Il settore dell'advertising digitale costituisce un dominio applicativo che impone l'elaborazione di elevati volumi di dati in tempo reale. La gestione delle campagne pubblicitarie necessita di un'infrastruttura di calcolo capace di scalare in risposta a traffico variabile e di fornire analisi relative al comportamento degli utenti. Per rispondere a tali requisiti operativi, i sistemi informativi moderni avviano transizioni verso architetture Cloud-Native, con l'obiettivo di superare i colli di bottiglia e i limiti di scalabilità delle soluzioni tradizionali.

Le infrastrutture legacy presentano spesso criticità operative dovute all'eterogeneità dei servizi e all'esecuzione parzialmente manuale dei processi di rilascio. L'assenza di un sistema di versionamento dichiarativo per l'infrastruttura rende complessa la tracciabilità delle modifiche e la riproducibilità degli ambienti operativi. Tale assenza comporta una scalabilità limitata e un'allocazione statica delle risorse computazionali, la quale risulta indipendente dal carico effettivo sostenuto dal sistema.

Questa tesi descrive la progettazione e l'implementazione di una nuova architettura finalizzata all'introduzione di standardizzazione e automazione per supportare carichi di lavoro intensivi. Per validare empiricamente il sistema, il progetto definisce un ambiente cloud di laboratorio configurato come setting simulato. L'infrastruttura elabora carichi di lavoro reali, interfacciandosi direttamente con le interfacce di programmazione (API) di provider esterni mediante l'impiego di parametri operativi e credenziali derivanti da un autentico caso industriale.

L'architettura proposta adotta il paradigma Serverless di Amazon Web Services (AWS) per astrarre la gestione dei server e utilizzare un modello di fatturazione legato all'uso effettivo delle risorse \cite{BerkeleyServerless, EivyEconomics}. Per l'automazione del provisioning, l'infrastruttura impiega Terraform, applicando il paradigma Infrastructure as Code (IaC) \cite{IaCGuerriero, TerraformUpAndRunning}. Ogni componente cloud risiede in repository Git \cite{Git}, consentendo la creazione di ambienti di sviluppo, staging e produzione riproducibili e integrando la gestione dell'infrastruttura nei flussi di revisione del software.

I processi di elaborazione dati seguono un pattern Event-Driven \cite{EventDrivenArchitecture}. Il flusso Extract, Transform, Load (ETL) si articola in microservizi asincroni. Un servizio di estrazione deposita i dati provenienti da molteplici sorgenti in un Data Lake basato su Amazon S3 \cite{S3Doc}. L'inserimento dei file innesca eventi che avviano le operazioni di trasformazione tramite AWS Lambda \cite{AWSLambdaDoc} o AWS Glue \cite{AWSGlueDoc}. Tali componenti applicano la logica di dominio e convertono i dati nel formato colonnare Parquet, ottimizzando i tempi di interrogazione tramite Amazon Athena \cite{Athena} e il consumo di risorse per le analisi di business intelligence.

Lo sviluppo dei servizi di backend sincroni adotta il framework Python FastAPI \cite{FastAPIDoc}. L'impiego di tale strumento agevola la validazione dei tipi a runtime, demandata alla libreria Pydantic, e supporta l'adozione di pattern architetturali volti a favorire modularità e testabilità del codice.

Un ulteriore aspetto trattato riguarda la standardizzazione dell'ambiente di sviluppo locale. Il progetto utilizza lo standard Docker \cite{DockerDoc} per il packaging delle funzioni Lambda, isolando le dipendenze e garantendo la parità di esecuzione tra l'ambiente di sviluppo e quello di produzione. Il ciclo di programmazione impiega container che emulano il runtime AWS e strumenti come LocalStack \cite{LocalStackPaper} per simulare localmente risorse cloud quali database DynamoDB, code SQS e bucket S3. Tale approccio abilita l'esecuzione di test d'integrazione sulla workstation, riducendo le latenze di rete e ottimizzando i tempi di iterazione.

Inoltre, i processi di Continuous Integration e Continuous Deployment (CI/CD) operano tramite GitHub Actions, applicando metodologie GitOps \cite{GitOps}. La pipeline gestisce il versionamento semantico, i changelog e il deployment delle risorse infrastrutturali, minimizzando l'intervento manuale e tracciando le modifiche \cite{AccelerateBook}.

Infine, per automatizzare la creazione di nuovi microservizi e standardizzare le architetture di base, la tesi presenta Service-Forge, un framework di automazione sviluppato ad hoc. Concepito nativamente come strumento a riga di comando (CLI) per l'esecuzione di operazioni di scaffolding, il sistema espone un'interfaccia basata sul Model Context Protocol (MCP) \cite{MCP}. Questa astrazione trasforma il framework in un server interrogabile in modo deterministico da agenti basati su Large Language Models (LLM) integrati negli ambienti di sviluppo (IDE). Il sistema traduce le specifiche di alto livello in impalcature di codice e pipeline CI/CD, assicurando che le nuove risorse soddisfino gli standard di sicurezza, e offre all'operatore la flessibilità di agire tramite terminale tradizionale o delegando l'orchestrazione all'assistente conversazionale.

Il presente documento articola i contenuti in cinque capitoli successivi a questa introduzione. La sequenza espositiva segue un percorso logico lineare: la trattazione parte dall'inquadramento teorico degli strumenti, attraversa l'analisi dei requisiti e la progettazione architetturale del sistema, per concludersi con la misurazione oggettiva delle prestazioni.

Il Capitolo \ref{chap:background} espone il background tecnologico. Questa porzione di testo fornisce i fondamenti teorici necessari alla comprensione dell'elaborato, analizzando i paradigmi Cloud-Native, l'architettura Serverless e le metodologie di gestione dell'infrastruttura come codice (IaC).

Il Capitolo \ref{chap:analisi} analizza il contesto operativo e definisce i requisiti del sistema. L'esposizione esamina i colli di bottiglia e le limitazioni strutturali delle architetture tradizionali, formalizzando i vincoli architetturali e operativi che guidano la riprogettazione dell'infrastruttura.

Il Capitolo \ref{chap:implementazione} descrive la progettazione e l'implementazione della soluzione proposta. Il testo illustra la topologia cloud adottata, il funzionamento delle pipeline asincrone per l'elaborazione dei dati, lo sviluppo dei servizi applicativi sincroni e l'integrazione del framework di scaffolding automatizzato.

Il Capitolo \ref{chap:caso_studio} quantifica le prestazioni dell'architettura tramite una validazione empirica. Il capitolo applica una metodologia di simulazione guidata da tracce per misurare i tempi di esecuzione, il consumo di risorse, l'efficienza algoritmica dei partizionamenti e l'abbattimento dei costi operativi.

Infine, il Capitolo \ref{chap:conclusioni} sintetizza il lavoro svolto. La conclusione confronta i risultati ottenuti con gli obiettivi iniziali e delinea le possibili evoluzioni future dell'infrastruttura.